\section{Introduction}
\label{sec:introduction}
Earth science seeks to derive testable scientific conclusions from observational data, from climate change and critical minerals to energy resources and disaster risk. Its research frontiers are data intensive, computationally demanding, and interdisciplinary, with increasing emphasis on predictive systems science driven by big data and artificial intelligence~\cite{zheng2025frontiers}. Progress requires organizing distributed data, tools, and methods into reproducible, traceable analytical workflows.

Such workflows typically require data retrieval, spatial or temporal alignment, variable derivation, statistical modeling, and result verification. Researchers manually connect these stages, despite differences in data formats and reporting standards. Scientific agents built around large language models (LLMs) offer a way to automate this process by translating natural-language goals into multistep plans and coordinating tools~\cite{feng2026earthagent,ding2025scitoolagent}. Crucially, later geoscientific operations depend on artifacts produced by earlier steps. Reliable tool orchestration is therefore essential to completing the analysis.

Existing geoscientific agents support multistep execution using frameworks such as ReAct~\cite{yao2023react}. GeoAgent~\cite{chen2024geoagent} and Earth-Agent~\cite{feng2026earthagent}, for example, use planning--execution--reflection architectures. However, their execution units are typically individual calls, individual tools, or entire Skills. Tool-level planning couples workflows to concrete interfaces, whereas whole-Skill reuse makes it difficult to extract and recombine local operations from different strategies. An unresolved problem is how to compose operations across strategies while preserving their prerequisites and usage constraints, and subsequently select suitable tool implementations.

Evaluation presents a related challenge. EarthSE assesses geoscientific reasoning through questions and exploratory dialogue without requiring actual tool execution~\cite{xu2025earthse}. GeoBenchX~\cite{krechetova2025geobenchx}, Earth-Bench~\cite{feng2026earthagent}, and GeoAgentBench~\cite{yu2026geoagentbench} examine multistep tool use; the latter two also assess processes and outcomes. Nevertheless, comprehensive evaluation across geoscientific research themes remains limited, particularly for distinguishing correct answers from complete analytical workflows and attributing orchestration failures to specific stages.

We introduce AgentGS and GeoToolEval to address these challenges. Our contributions are:
\begin{itemize}
\item AgentGS separates Skills, Capabilities, and Tools and completes missing capabilities through a hypergraph. This supports operation reuse across strategies, constraint preservation, and executable tool selection.
\item GeoToolEval comprises 603 evaluation records across ten geoscientific research-frontier themes. Single-step and multistep tasks support joint analysis of tool dependencies, composition, intermediate-artifact use, and final answers.
\item AgentGS achieves a mean Final Answer Score (FAS) of 0.8731 on GeoToolEval, outperforming the evaluated models and agents. Ablations identify the contributions of hierarchical planning, predecessor expansion, capability completion, and execution recovery.
\end{itemize}
